% G. H. Hardy, Divergent Series (Oxford, Clarendon Press, 1949), the start
% of chapter III, pp. 42-43: theorems, fraktur letters, a footnote.

\pagenumber{42}
\chapter[III]{General theorems}

\section{Generalities concerning linear transformations.} The theory of
divergent series is concerned with generalizations of the notion of the
limit of a sequence $(s_n)$, which are usually effected by an auxiliary
sequence of linear means of $s_n$. Thus in \S1.3 we defined the (C,~1)
limit of $(s_n)$, or the (C,~1) sum of $\sum a_n$, as the limit of
$$\no\label{c1} t_m = \frac{s_0+s_1+\dots+s_m}{m+1}$$
when $m\to\infty$; and the A limit of $(s_n)$, or the A sum of $\sum a_n$,
as the limit of
$$\no\label{abel} t(x) = \sum a_nx^n = \sum x^n(1-x)s_n$$
when $x\to1$ through values less than 1. In each case the auxiliary means
are of the form
$$\no\label{tm} t_m = \sum c_{m,n}s_n \quad (m=0,1,2,\dots)$$
or
$$\no\label{tx} t(x) = \sum c_n(x)s_n,$$
where $x$ is a continuous parameter.\footnote{Summations are over
$0,1,2,\dots$ when there is no indication to the contrary. The variable of
summation will not be shown explicitly unless this is necessary to avoid
ambiguity: it is obvious, for example, that in \eqref{tm} and \eqref{tx}
the summation must be with respect to $n$, so that, for example,
$\sum c_{m,n}s_n$ means $\sum_{n=0}^\infty c_{m,n}s_n$.} Thus in \eqref{c1}
$$c_{m,n} = (m+1)^{-1} \quad (0\le n\le m), \qquad c_{m,n} = 0 \quad (n>m),$$
and in \eqref{abel} $c_n(x)=x^n(1-x)$. In one case they depend on an
integral parameter $m$, in the other on a continuous parameter $x$; but, as
we shall see, this distinction is not very important. For the moment we
consider means of the type \eqref{tm}.

We call the system of equations \eqref{tm}, which we may write shortly as
$$\no t = \mathrm{T}(s),$$
a \emph{linear transformation} T; $t_m$ the \emph{transform} of $s_n$ by T;
and the matrix
$$|\mathrm{T}| = (c_{m,n}),$$
in which $c_{m,n}$ is the element in the $m$th row and $n$th column, the
\emph{matrix of}~T.

\section{Regular transformations.} The most important transformations are
\emph{regular}. We say that T is regular if
$$\no t_m\to s \quad (m\to\infty)$$
whenever
$$\no s_n\to s \quad (n\to\infty).$$
We regard the first assertion as including that of the existence of $t_m$
for every $m$, i.e. the convergence of all the series \eqref{tm}. Thus,
after Cauchy's theorem quoted in \S1.4, the transformation \eqref{c1} is
regular.

There is an important theorem, due to Toeplitz and Schur, which states
necessary and sufficient conditions for the regularity of T. We prove this
theorem (Theorem~2) in \S3.3: it is convenient to associate it with two
other theorems of a similar character concerning different classes of
transformations. We call the class of linear transformations
$\mathfrak{T}$, the class of regular transformations $\mathfrak{T}_r$. The
class $\mathfrak{T}_c$ is the class of transformations which transform all
convergent sequences into convergent sequences, i.e. transformations such
that the convergence of $s_n$ to $s$ implies the convergence of $t_m$ to
\emph{some} limit $t$. Thus $\mathfrak{T}_r$ is the subclass of
$\mathfrak{T}_c$ in which $t$ is necessarily the same as $s$. The class
$\mathfrak{T}_c^*$ is the class of transformations which convert all
bounded sequences into convergent sequences, i.e. transformations such that
$s_n=O(1)$ implies $t_m\to t$. It is plain that $\mathfrak{T}_c^*$ is also
a subclass of $\mathfrak{T}_c$; but $\mathfrak{T}_c^*$ and $\mathfrak{T}_r$
are, as we shall see, mutually exclusive. We shall prove the following
three theorems.

\begin{theorem}\label{thm-c} In order that \textup{T} should belong to
$\mathfrak{T}_c$, it is necessary and sufficient \textup{(i)} that
$$\no \gamma_m = \sum|c_{m,n}| < H,$$
where $H$ is independent of $m$; \textup{(ii)} that
$$\no c_{m,n}\to\delta_n,$$
for each $n$, when $m\to\infty$; \textup{(iii)} that
$$\no c_m = \sum c_{m,n}\to\delta$$
when $m\to\infty$. In these circumstances $\sum\delta_n$ is absolutely
convergent, and
$$\no t_m\to t = \delta s+\sum\delta_n(s_n-s) = s(\delta-\sum\delta_n)+\sum\delta_ns_n,$$
when $m\to\infty$, whenever $s_n\to s$.
\end{theorem}

Here, of course, the limits $\delta_n$ and $\delta$ are finite.

\begin{theorem} In order that \textup{T} should belong to $\mathfrak{T}_r$
\textup{(}i.e. that \textup{T} should be regular\textup{)}, it is necessary
and sufficient that the conditions of Theorem~\ref{thm-c} should be
satisfied, that $\delta_n=0$ for each $n$, and that $\delta=1$.
\end{theorem}
